\chapter{Bose--Einstein Distribution}

\section*{Introduction}

Imagine a crowd of identical people who prefer to cluster together rather than spread apart. The Bose--Einstein distribution quantifies this clustering tendency. Each quantum state is like a seat in a theater: the more people already sitting there, the more attractive that seat becomes to the remaining crowd. This ``bosonic stimulation'' is not a psychological effect but a consequence of quantum indistinguishability: bosons are identical particles that cannot be told apart, so swapping two bosons in the same state does not produce a new configuration---it amplifies the existing one. The result is that bosons overwhelmingly prefer to share states, in sharp contrast to fermions, which obey Pauli exclusion and refuse to occupy the same state.


The mean occupation number of a single-particle state with energy $E$ is:
\[
\langle n \rangle = \frac{1}{e^{(E - \mu)/kT} - 1}
\]
where $k$ is Boltzmann's constant, $T$ is the absolute temperature, and $\mu$ is the chemical potential.

\section*{Fundamental Equations Used}

The derivation starts from the grand canonical ensemble, with a factorised grand partition function over single-particle quantum states.

\section*{Assumptions}

\textbf{Assumptions:} The particles are identical bosons (integer spin: 0, 1, 2, \ldots). The system is in thermal equilibrium. The single-particle energy levels are well-defined.


\textbf{Limitations:} Does not describe interactions between bosons (ideal Bose gas). Breaks down for fermions (which require the Fermi--Dirac distribution with a ``$+1$'' in the denominator). At very high densities or strong interactions, the mean-field description becomes inaccurate.


\section*{Mathematical Derivation}

\textbf{Step 1: Define the grand canonical ensemble.}

Consider a system of non-interacting identical bosons in thermal equilibrium with a heat bath at temperature $T$ and chemical potential $\mu$. Each single-particle quantum state $k$ with energy $E_k$ is an independent subsystem that can hold any number of bosons $n_k = 0, 1, 2, \ldots$ The grand partition function factorizes over all single-particle states:
\begin{equation}
\mathcal{Z} = \prod_k \mathcal{Z}_k, \qquad \mathcal{Z}_k = \sum_{n_k=0}^{\infty} e^{-\beta(E_k - \mu)n_k}
\end{equation}
where $\beta = 1/(kT)$.

\textbf{Step 2: Evaluate the single-state partition function.}

Since $e^{-\beta(E_k - \mu)}$ is a number less than 1 (for $\mu < E_k$, which is required for convergence), the geometric series sums exactly:
\begin{equation}
\mathcal{Z}_k = \sum_{n_k=0}^{\infty} \left[e^{-\beta(E_k - \mu)}\right]^{n_k} = \frac{1}{1 - e^{-\beta(E_k - \mu)}}
\end{equation}

\textbf{Step 3: Compute the grand potential.}

The logarithm of the grand partition function gives the grand potential:
\begin{equation}
\ln \mathcal{Z} = -\sum_k \ln\left(1 - e^{-\beta(E_k - \mu)}\right)
\end{equation}

\textbf{Step 4: Derive the mean occupation number.}

The average number of particles in state $k$ is:
\begin{align}
\langle n_k \rangle &= \frac{1}{\beta} \frac{\partial \ln \mathcal{Z}_k}{\partial \mu} \notag \\[6pt]
&= \frac{1}{\beta} \frac{\partial}{\partial \mu} \left[-\ln\left(1 - e^{-\beta(E_k - \mu)}\right)\right] \notag \\[6pt]
&= \frac{1}{\beta} \cdot \frac{\beta \, e^{-\beta(E_k - \mu)}}{1 - e^{-\beta(E_k - \mu)}} \notag \\[6pt]
&= \frac{e^{-\beta(E_k - \mu)}}{1 - e^{-\beta(E_k - \mu)}}
\end{align}

\textbf{Step 5: Simplify to the standard form.}

Multiply numerator and denominator by $e^{\beta(E_k - \mu)}$:
\begin{equation}
\boxed{\langle n_k \rangle = \frac{1}{e^{\beta(E_k - \mu)} - 1} = \frac{1}{e^{(E_k - \mu)/kT} - 1}}
\end{equation}
This is the Bose--Einstein distribution. The $-1$ in the denominator is the hallmark of bosonic statistics, arising from the symmetric (bosonic) commutation relations. For fermions, the Pauli exclusion principle restricts $n_k$ to 0 or 1, and a $+1$ appears instead (Fermi--Dirac distribution).

\textbf{Step 6: Verify limiting cases.}

\emph{Classical limit:} When $\langle n_k \rangle \ll 1$, we have $e^{\beta(E_k - \mu)} \gg 1$, so the $-1$ is negligible and $\langle n_k \rangle \approx e^{-\beta(E_k - \mu)}$, which is the Maxwell--Boltzmann distribution.

\emph{Photon gas ($\mu = 0$):} For photons, particle number is not conserved, so $\mu = 0$. The distribution becomes $\langle n_k \rangle = 1/(e^{\beta E_k} - 1)$, the Planck distribution for blackbody radiation.

\section*{Characteristics of the Equation}

\begin{itemize}
\item \textbf{Bosonic stimulation:} The factor $-1$ in the denominator means occupation grows faster than linearly with decreasing energy---each particle in a state makes it more probable that another will join.
\item \textbf{Classical limit:} When $\langle n \rangle \ll 1$, the distribution reduces to the Maxwell--Boltzmann form $\langle n \rangle \approx e^{-(E-\mu)/kT}$.
\item \textbf{Condensation:} When $T$ drops below a critical temperature $T_c$, the ground state becomes macroscopically occupied---a Bose--Einstein condensate forms.
\item \textbf{Photon statistics:} For photons, $\mu = 0$ (photon number is not conserved), giving the Planck distribution $n = 1/(e^{E/kT} - 1)$.
\item \textbf{Exchange symmetry:} The distribution arises from the requirement that the many-body wave function is symmetric under interchange of any two bosons.
\end{itemize}
\section*{Solution Methods}

\begin{enumerate}
\item \textbf{Grand canonical ensemble:} Fix $T$, $V$, and $\mu$. The grand partition function factorizes: $\mathcal{Z} = \prod_{\mathbf{k}} (1 - e^{-(E_{\mathbf{k}}-\mu)/kT})^{-1}$. Differentiating $\ln \mathcal{Z}$ with respect to $\mu$ gives $\langle n \rangle$.
\item \textbf{Semiclassical density of states:} For free particles in 3D, $\rho(E) = \frac{V}{4\pi^2}\left(\frac{2m}{\hbar^2}\right)^{3/2} \sqrt{E}$. Integrating $\langle n \rangle \, \rho(E) \, dE$ gives the total particle number; matching to $N$ determines $\mu(T)$.
\item \textbf{Numerical methods:} For lattice models or interacting systems, quantum Monte Carlo or dynamical mean-field theory (DMFT) can be used beyond the ideal gas approximation.
\end{enumerate}
\section*{Verifications}

\begin{itemize}
\item \textbf{Planck spectrum fit:} The measured blackbody spectrum matches the Bose--Einstein distribution with $\mu = 0$ across all wavelengths and temperatures to extraordinary precision.
\item \textbf{BEC onset:} In cold-atom experiments, BEC formation is observed as a sharp peak in the velocity distribution at $T < T_c$, matching the ideal-gas prediction.
\item \textbf{Specific heat anomaly:} The $\lambda$-transition in liquid helium at $T_\lambda = 2.17$\,K is a signature of condensation, modified by interactions.
\item \textbf{Photon bunching:} Measured photon statistics from thermal light show super-Poissonian bunching consistent with Bose--Einstein statistics.
\end{itemize}
\section*{Applications}

\begin{itemize}
\item \textbf{Blackbody radiation:} Setting $\mu = 0$ for photons yields the Planck spectrum, from which the Stefan--Boltzmann law and Wien's displacement law follow.
\item \textbf{Bose--Einstein condensation:} Cooling dilute alkali gases to nanokelvin temperatures produces macroscopic condensates used to study superfluidity, quantum vortices, and atom lasers.
\item \textbf{Phonon heat capacity:} The Bose--Einstein distribution for phonons explains why the heat capacity of solids drops as $T^3$ at low temperatures (Debye model).
\item \textbf{Superfluidity:} Liquid $^4$He becomes a superfluid below 2.17\,K---a macroscopic quantum phenomenon directly linked to Bose--Einstein statistics.
\item \textbf{Lasers:} Stimulated emission---the process that makes lasers work---is the electromagnetic analogue of bosonic stimulation.
\end{itemize}
\section*{What Richard Feynman Would Have Asked}

\subsection*{1. Plain-English Translation}
\textit{``If you had to explain this equation to someone who never took physics, what would you say?''}

The Bose--Einstein distribution says: ``Particles that are truly identical prefer to be together.'' If you have a party where everyone is wearing the same mask---indistinguishable---people naturally cluster, because there is no social anxiety about being different. The equation gives the exact probability that a given energy state has a certain number of these identical particles in it. When the temperature is low enough, a huge fraction of all particles collapses into the lowest energy state---like everyone in the theater rushing to the front row because the view must be spectacular.

The minus sign in the denominator is the mathematical expression of this preference. Compare it to the Fermi--Dirac distribution, which has a $+1$: fermions are like people who insist on having their own row. Bosons are like people who are happy to share---and the more people already in a row, the more who want to join.

\subsection*{2. Localized Mechanics}
\textit{``What is happening at the quantum level when a particle joins a state that already has particles in it?''}

At the level of individual quantum events, bosonic stimulation is not a force or an interaction---it is a consequence of quantum indistinguishability. When two identical bosons are in the same state, the wave function is symmetric under their exchange: swapping them does not produce a new quantum state. This means there is only one possible configuration for having two particles in one state, versus multiple configurations for distributing particles across different states. The counting is fundamentally different from the classical case.

He would press you on the mechanism: the particle does not ``know'' that other particles are present. Rather, the quantum mechanical counting of states---the symmetry of the wave function---makes the final state with two particles in one level more probable than the state with one particle in each of two levels. The stimulation factor $(n+1)$ for adding one more particle to a state with $n$ particles comes directly from this counting.

\subsection*{3. Testing Extreme Limits}
\textit{``What happens at absurdly high or low temperatures?''}

At extremely high temperatures ($kT \gg E$ for all states), the exponential becomes very large, the $-1$ becomes negligible, and the Bose--Einstein distribution reduces to the classical Maxwell--Boltzmann distribution. Quantum statistics become irrelevant---this is why classical thermodynamics works at room temperature for most gases.

At temperatures approaching absolute zero, $\mu$ approaches the ground state energy $E_0$, and the ground state occupation diverges---the condensate forms. In a real system, interactions eventually stabilize the condensate. The critical temperature for BEC in a 3D ideal gas is $T_c = \frac{2\pi\hbar^2}{mk}\left(\frac{n}{\zeta(3/2)}\right)^{2/3}$, where $\zeta(3/2) \approx 2.612$ is the Riemann zeta function. Below this, the condensate fraction is $1 - (T/T_c)^{3/2}$.

\subsection*{4. Dimensionality and Geometry}
\textit{``How does the dimensionality of space affect Bose--Einstein condensation?''}

In 3D, the density of states $\rho(E) \propto \sqrt{E}$ grows with energy, allowing a finite fraction of particles to remain in excited states at low temperature. BEC occurs at a finite critical temperature. In 2D, $\rho(E) = \text{const}$---the density of states is flat, and the integral diverges logarithmically. Strictly speaking, BEC does not occur in a uniform 2D system at finite temperature (related to the Hohenberg--Mermin--Wagner theorem). However, in a 2D harmonic trap, a condensate can form because the trap modifies the effective density of states. In 1D, $\rho(E) \propto 1/\sqrt{E}$, making BEC even harder to achieve.

\subsection*{5. Cross-Disciplinary Connection}
\textit{``Where else does this exact same mathematical structure appear outside of quantum physics?''}

The Bose--Einstein distribution is a special case of the statistics governing any system of indistinguishable objects that can share states. In network theory, the ``rich get richer'' phenomenon---where nodes with many connections tend to acquire even more---is mathematically analogous to bosonic stimulation: the probability of acquiring a new connection is proportional to $(n + 1)$. This produces scale-free networks with power-law degree distributions, similar to the tails of the Bose--Einstein distribution at low energies.

In economics, the same structure appears in models of wealth distribution where identical financial instruments are shared among agents, and in queuing theory where indistinguishable customers join servers. The mathematical universality of the Bose--Einstein distribution---rooted in the symmetry of indistinguishable entities sharing states---means it appears far more broadly than its original quantum context suggests.


%=============================================================
% Chapter 14: Boundary Layer Equation
%=============================================================
\section*{FAQ}

\textbf{Q1: Why does the minus sign in the denominator matter so much?}

A: The minus sign (versus the $+1$ in Fermi--Dirac) is the mathematical signature of bosonic statistics. It means the occupation number can become arbitrarily large---there is no exclusion principle. When $E = \mu$, the denominator vanishes and $\langle n \rangle$ diverges: this is condensation.

\textbf{Q2: Can fermions form a condensate?}

A: Not directly, because of Pauli exclusion. However, two fermions can form a bound pair (a composite boson), and these pairs can condense---this is the mechanism behind superconductivity (Cooper pairs) and superfluidity in $^3$He.

\textbf{Q3: Why is the chemical potential negative for an ideal Bose gas?}

A: To keep all occupation numbers positive and finite, we need $\mu < E_0$ for all energy levels. For a free gas, $E_0 = 0$, so $\mu < 0$. As $T$ decreases, $\mu$ approaches zero from below, and at $T_c$ it reaches zero---below $T_c$, the ground state is macroscopically occupied.
\section*{Important Bibliography}

\begin{enumerate}
\item Pathria, R.K. and Beale, P.D., \textit{Statistical Mechanics}, 4th ed. (Academic Press, 2022), Chapter 7.
\item Bose, S.N., ``Plancks Gesetz und Lichtquantenhypothese,'' \textit{Zeitschrift f\"ur Physik} 26, 178--181 (1924).
\item Einstein, A., ``Quantentheorie des einatomigen idealen Gases,'' \textit{Sitzungsberichte der Preussischen Akademie der Wissenschaften}, 261--267 (1925).
\item Pethick, C.J. and Smith, H., \textit{Bose--Einstein Condensation in Dilute Gases}, 2nd ed. (Cambridge University Press, 2008), Chapter 1.
\item Mandl, F., \textit{Statistical Physics}, 3rd ed. (Wiley, 1988), Chapter 8.
\end{enumerate}

